\documentclass[border=5mm]{standalone}
% ==============================================================================
% SETUP.TEX - CONFIGURACIÓN GLOBAL DEL PROYECTO
% ==============================================================================
% Este archivo centraliza todos los paquetes y estilos.
% Debe incluirse en el preámbulo del libro principal y de los archivos standalone.
% \PassOptionsToPackage{gray}{xcolor}
% --- 1. CODIFICACIÓN E IDIOMA ---
% fix-cm habilita las fuentes Computer Modern en tamaños arbitrarios (por debajo
% de 5 pt, entre otros). Debe cargarse antes que fontenc.
\usepackage{fix-cm}
\usepackage[utf8]{inputenc}
\usepackage[T1]{fontenc}
\usepackage[spanish, es-tabla]{babel}  
\usepackage{csquotes} % Recomendado para babel y citas

\usepackage{import}
% mode=image: Usa el PDF pre-compilado, nunca intenta recompilar.
% Debes compilar cada figura manualmente antes de compilar el libro.
\usepackage[mode=image, subpreambles=false]{standalone}

% --- 2. MATEMÁTICAS Y CIENCIAS ---
\usepackage{amsmath, amssymb, amsfonts, mathtools}
\usepackage{enumitem}

% LaTeX trae \DeclareMathSizes{5}{5}{5}{5}, así que a 5 pt (\tiny) los subíndices
% salen del mismo tamaño que la base: en las etiquetas de figuras el M de
% $\omega_M$ queda desproporcionado. Se restablece la proporción habitual.
\DeclareMathSizes{5}{5}{3.7}{3}

% --- 3. DISEÑO Y GEOMETRÍA --- 
% Unificamos los márgenes para todo el documento
\usepackage[left=2.5cm,right=2.5cm,top=3cm,bottom=3cm]{geometry}
\makeatletter
\ifstandalone % Evita que geometry dañe el recorte de las figuras bajándolas 3cm
  \@ifclasswith{standalone}{crop=false}{
    % Es un capítulo/sección: Dejamos que geometry actúe.
  }{
    % Es una figura: Neutralizamos geometry para que no las recorte mal.
    \geometry{pass}
  }
\fi
\makeatother
\usepackage{parskip} % Espaciado entre párrafos sin sangría
\usepackage{float}   % Para usar [H] en figura

% Ajustes de flotantes para evitar espacios verticales exagerados
\renewcommand{\topfraction}{0.95}      % Permite que las figuras ocupen hasta el 95% superior
\renewcommand{\bottomfraction}{0.95}   % Permite que las figuras ocupen hasta el 95% inferior
\renewcommand{\textfraction}{0.05}     % Permite hojas con tan solo un 5% de texto
\renewcommand{\floatpagefraction}{0.8} % Requiere que una página de solo figuras esté 80% llena
\setcounter{topnumber}{4}
\setcounter{bottomnumber}{4}
\setcounter{totalnumber}{10}

% --- 4. GRÁFICOS Y TABLAS ---
\usepackage{graphicx}
\usepackage{booktabs} % Tablas profesionales
\usepackage{caption}  % Control de captions y \captionof
\usepackage{tikz}
\usetikzlibrary{arrows.meta, calc, angles, quotes, babel, backgrounds}
\usepackage{pgfplots}
\usepgfplotslibrary{groupplots}
\usepgfplotslibrary{external}
\usepgfplotslibrary{patchplots}
\pgfplotsset{compat=1.18}
\usepackage[american, straightvoltages]{circuitikz}

% --- 5. COLORES Y CAJAS ---

\usepackage[table]{xcolor}
\usepackage{tcolorbox}

% Definición de colores corporativos/estilo
\definecolor{utnblue}{RGB}{0, 51, 102}
\definecolor{codegreen}{rgb}{0,0.6,0}
\definecolor{codegray}{rgb}{0.5,0.5,0.5}
\definecolor{codepurple}{rgb}{0.58,0,0.82}
\definecolor{backcolour}{rgb}{0.96,0.96,0.96}

% --- 6. ENTORNOS PERSONALIZADOS (CAJAS) ---

% Caja "Resultado" (Azul Corporativo) — Definiciones, fórmulas clave, resultados importantes.
% Uso: \begin{resultado}[Fórmula de Euler] ... \end{resultado}
\newtcolorbox{resultado}[1][]{
    colback=utnblue!5!white,
    colframe=utnblue,
    title=\textbf{#1},
    fonttitle=\bfseries,
    sharp corners,
    boxrule=0.5mm
}

% Caja "Nota" (Gris) — Advertencias, aplicaciones, contexto secundario.
% Uso: \begin{nota}[Cuidado con la calculadora] ... \end{nota}
% Uso: \begin{nota} ... \end{nota}  → título por defecto "Nota"
\newtcolorbox{nota}[1][Nota]{
    colback=codegray!5!white,
    colframe=codegray,
    title=\textbf{#1},
    fonttitle=\bfseries,
    sharp corners,
    boxrule=0.5mm
}

% Caja inline para resaltar un valor y enlazarlo con su otra aparición en una tabla
% o en una cuenta: mismo color, mismo valor.
% Uso: \marcavalor{$4$}  o  \marcavalor[codegreen]{$4$}
\newtcbox{\marcavalor}[1][utnblue]{
    on line,
    boxsep=1pt, left=2pt, right=2pt, top=1pt, bottom=1pt,
    colback=#1!10!white,
    colframe=#1,
    boxrule=0.4pt,
    arc=2pt
}

% --- 7. CÓDIGO FUENTE (LISTINGS) ---
\usepackage{listings}

\lstdefinestyle{miestilo}{
    backgroundcolor=\color{backcolour},   
    commentstyle=\color{codegreen},
    keywordstyle=\color{magenta},
    numberstyle=\tiny\color{codegray},
    stringstyle=\color{codepurple},
    basicstyle=\ttfamily\footnotesize,
    breakatwhitespace=false,         
    breaklines=true,                 
    captionpos=b,                    
    keepspaces=true,                 
    numbers=left,                    
    numbersep=5pt,                  
    showspaces=false,                
    showstringspaces=false,
    showtabs=false,                  
    tabsize=2,
    frame=single,                    
    rulecolor=\color{codegray}
}

\lstset{style=miestilo}
% Soporte para tildes y ñ en bloques de código
\lstset{literate={ñ}{{\~n}}1 {á}{{\'a}}1 {é}{{\'e}}1 {í}{{\'i}}1 {ó}{{\'o}}1 {ú}{{\'u}}1}

% Operadores matemáticos personalizados

% Vector en sentido discreto (lista finita de coordenadas): negrita + flecha.
% Uso: \vect{v}, \vect{e}_k, \vect{0}.
\newcommand{\vect}[1]{\vec{\mathbf{#1}}}

% Fecha de versión y comando para capítulos standalone
\providecommand{\verdate}{\the\year.\ifnum\month<10 0\fi\the\month.\ifnum\day<10 0\fi\the\day}
\newcommand{\chapterversion}{%
    \vspace{-1.5cm}\begin{flushright}\small\textit{\color{codegray}Versión~\verdate}\end{flushright}\vspace{0.5cm}%
}

% --- 8. HIPERVÍNCULOS (DEBE SER EL ÚLTIMO PAQUETE) ---
\usepackage[colorlinks=true, linkcolor=utnblue, urlcolor=utnblue, citecolor=utnblue]{hyperref}

% --- 9. REFERENCIAS CRUZADAS ENTRE CAPÍTULOS STANDALONE ---
% Cuando se compila un capítulo standalone, xr-hyper lee los .aux de los
% otros capítulos (si existen) para resolver \ref{} a labels externos.
% En la compilación del libro completo este bloque no se activa.
\ifstandalone
  \usepackage{xr-hyper}
  \makeatletter
  \newcommand{\xrchap}[1]{%
    \edef\xrc@self{\detokenize{Capitulo#1}}%
    \edef\xrc@job{\jobname}%
    \ifx\xrc@self\xrc@job\else
      \IfFileExists{../#1capitulo/Capitulo#1.aux}%
        {\externaldocument{../#1capitulo/Capitulo#1}}{}%
    \fi
  }
  \makeatother
  \xrchap{01}\xrchap{02}\xrchap{03}\xrchap{04}
  \xrchap{05}\xrchap{06}\xrchap{07}\xrchap{08}
  \xrchap{09}\xrchap{10}\xrchap{11}\xrchap{12}
  \xrchap{13}
\fi
% \PassOptionsToPackage{gray}{xcolor}
\usetikzlibrary{decorations.pathmorphing, patterns}

\begin{document}
    \begin{tikzpicture}[>=Latex, font=\small,
        spring/.style={decorate, decoration={zigzag, segment length=5, amplitude=3, pre length=3, post length=3}},
    ]
        % =============================================
        % PANEL IZQUIERDO: SISTEMA MECÁNICO
        % =============================================
        \node[font=\bfseries, anchor=south] at (2, 5.5) {Sistema Mecánico};

        % Techo
        \fill[pattern=north east lines, pattern color=gray] (0, 5) rectangle (4, 5.2);
        \draw[very thick] (0, 5) -- (4, 5);

        % Resorte
        \draw[spring, thick] (1, 5) -- (1, 3) node[midway, left=3mm, font=\footnotesize] {$k$};

        % Amortiguador simplificado
        \draw[thick] (3, 5) -- (3, 4.2);
        \draw[thick, fill=gray!15] (2.7, 4.2) rectangle (3.3, 4.0);
        \draw[thick] (3, 4.0) -- (3, 3);
        \draw[thick] (2.7, 4.2) -- (2.7, 3.8);
        \draw[thick] (3.3, 4.2) -- (3.3, 3.8);
        \node[right=3mm, font=\footnotesize] at (3.4, 4.1) {$b$};

        % Masa
        \draw[thick, fill=utnblue!15] (0.3, 2.2) rectangle (3.7, 3);
        \node at (2, 2.6) {\textbf{$m$}};

        % Fuerza
        \draw[->, ultra thick, red!70!black] (2, 2.2) -- (2, 1) node[right, font=\footnotesize] {$f(t)$};

        % Desplazamiento
        \draw[<->, thick, codegreen!70!black] (4.3, 2.6) -- (4.3, 1.3) node[midway, right, font=\footnotesize] {$y(t)$};

        % Ecuación (entrada roja, salida verde)
        \node[draw, thick, rounded corners, fill=orange!8, inner sep=6pt] at (2, -0.3) {
            $m\ddot{\textcolor{codegreen!70!black}{y}} + b\dot{\textcolor{codegreen!70!black}{y}} + k\,\textcolor{codegreen!70!black}{y} = \textcolor{red!70!black}{f(t)}$
        };

        % Tabla de correspondencias
        \node[font=\footnotesize, anchor=north] at (2, -1.2) {
            \begin{tabular}{rl}
            Masa: & $m$ \\
            Fricción: & $b$ \\
            Rigidez: & $k$ \\
            Fuerza (entrada): & $\textcolor{red!70!black}{f(t)}$ \\
            Posición (salida): & $\textcolor{codegreen!70!black}{y(t)}$
            \end{tabular}
        };

        % =============================================
        % SIGNO DE EQUIVALENCIA
        % =============================================
        \node[font=\Huge, gray] at (6.5, 2.5) {$\equiv$};
        \node[font=\footnotesize\itshape, gray, text width=2.5cm, align=center] at (6.5, 1) {Mismo modelo,\\distinto fenómeno};

        % =============================================
        % PANEL DERECHO: CIRCUITO RLC (circuitikz)
        % =============================================
        \node[font=\bfseries, anchor=south] at (11, 5.5) {Sistema Eléctrico};

        % Circuito RLC serie con circuitikz (componentes más compactos)
        \ctikzset{bipoles/length=0.8cm}
        \draw[thick] (9, 1.5) 
            to[sV, l_=$\textcolor{red!70!black}{v(t)}$] (9, 5)  % Fuente de tensión (entrada)
            -- (13, 5)                                            % Cable superior
            to[R, l_=$R$] (13, 3.75)                              % Resistencia
            to[L, l_=$L$] (13, 2.5)                               % Inductancia
            to[C, l_=$C$] (13, 1.5)                               % Capacitor
            -- (9, 1.5);                                           % Cable inferior

        % Flecha de corriente (salida)
        \draw[->, thick, black] (10., 4.7) -- (12, 4.7) 
            node[midway, below, font=\footnotesize] {$i(t)$};

        % Ecuación (entrada roja, salida verde)
        \node[draw, thick, rounded corners, fill=orange!8, inner sep=6pt] at (11, -0.3) {
            $L\ddot{\textcolor{codegreen!70!black}{q}} + R\dot{\textcolor{codegreen!70!black}{q}} + \frac{1}{C}\textcolor{codegreen!70!black}{q} = \textcolor{red!70!black}{v(t)}$
        };

        % Tabla de correspondencias
        \node[font=\footnotesize, anchor=north] at (11, -1.2) {
            \begin{tabular}{rl}
            Inductancia: & $L$ \\
            Resistencia: & $R$ \\
            Elastancia: & $1/C$ \\
            Tensión (entrada): & $\textcolor{red!70!black}{v(t)}$ \\
            Carga (salida): & $\textcolor{codegreen!70!black}{q(t)}$
            \end{tabular}
        };

    \end{tikzpicture}
\end{document}
